% ===============================================================
%  Web プログラミング レポート
%  コンパイル: lualatex report.tex （TeX Live 2025 / LuaLaTeX）
% ===============================================================
\documentclass[a4paper,11pt]{ltjsarticle}

\usepackage{listings}
\usepackage{xcolor}
\usepackage{url}

% --- ソースコード表示の設定（日本語コメントを含む JavaScript 用）---
\definecolor{codegray}{gray}{0.4}
\definecolor{codekeyword}{RGB}{0,0,180}
\definecolor{codestring}{RGB}{160,0,0}

% listings は JavaScript を標準サポートしないため，コメント・文字列・キーワードを手動定義する
\lstset{
  basicstyle=\ttfamily\small,
  commentstyle=\color{codegray},
  keywordstyle=\color{codekeyword},
  stringstyle=\color{codestring},
  morecomment=[l]{//},
  morecomment=[s]{/*}{*/},
  morestring=[b]",
  morestring=[b]',
  morekeywords={const,let,var,function,return,require,if,else,for,switch,case,break,default,new},
  numbers=left,
  numberstyle=\tiny\color{codegray},
  stepnumber=1,
  numbersep=8pt,
  frame=single,
  breaklines=true,
  breakindent=20pt,
  showstringspaces=false,
  tabsize=2,
  columns=fullflexible,
  keepspaces=true,
}

\title{Web アプリの作成と Azure 上での公開}
\author{学籍番号: 25G1104 \quad 氏名: 沼田 晴}
\date{\today}

\begin{document}
\maketitle

\section{目的}

本システムには、残念なポケモン図鑑、最強剣一覧、特殊羽一覧の3つの機能が統合されている．
残念なポケモン図鑑では，ポケモンのコンテンツでは多く存在する図鑑説明が辛辣なポケモンを集め，
ユーザに提供することで，ユーザに楽しんでもらうことを目的としている．最強剣一覧システムでは，チョコットランドに登場する最強剣の情報を集約し，ユーザに提供すること
で，ユーザがゲーム内での装備選択，作成を行う際の参考情報を提供することを目的としている．特殊羽一覧システムでは，チョコットランドに登場する特殊羽根装備の情報を集約し，ユーザに提供する
ことで，ユーザがゲーム内での装備の購入をする際やドーピングを行う際の参考情報を提供するこ
とを目的としている．

\section{公開URL}

本システムは Azure 上の仮想マシンで公開しており，以下の URL からアクセスできる．

\begin{itemize}
  \item トップページ（一覧メニュー）: \url{http://20.214.236.18:8080/}
\end{itemize}

\section{サーバー側プログラム}

サーバー側のプログラムを以下に示す．

\lstinputlisting{report-app5.js}

\end{document}
